\chapter{Introduction}
\label{ch:intro}

This chapter introduces the problem of memory design for LLM-based trading agents, explains the idea behind our self-correcting agent memory, states our objectives and methodology, and summarises the outcomes and the organisation of the report.

\section{Project Overview}
Large language models (LLMs) can read news, filings and analyst notes and turn them into trading decisions with a written rationale. Because a model's context window is limited and markets move every day, such agents need a \emph{memory}: a store of past information from which a few relevant items are retrieved for each decision. The most widely cited design, FinMem~\cite{yu2023finmem}, splits long-term memory into shallow, intermediate and deep layers for daily news, quarterly (10-Q) and annual (10-K) reports, each with its own decay speed, ranks memories by recency, relevancy and importance, promotes useful memories to deeper layers, and reports large gains over buy-and-hold and reinforcement-learning agents on five U.S. stocks.

Two problems motivate this project. First, \textbf{memory design}: such memories treat all sources within a layer as equally trustworthy, keep facts that later news contradicts until they decay, store the same story many times when it is re-reported, sample importance at random, ignore related companies, and do not learn from the realised outcomes of their own decisions. An agent memory should instead \emph{remember} what is relevant now, \emph{verify} it against source reliability, contradicting news and past outcomes, and only then let the agent \emph{act}. Our memory is built around this idea: four components improve what is \emph{remembered} (duplicate consolidation, a situational query, data-driven importance and related-company news), three \emph{verify} it (learned source credibility, contradiction retirement and a check of how acting in similar past situations turned out), and the agent then \emph{acts} through a decision policy and, in the application, through hard budget guardrails. Second, \textbf{evaluation}: published results for LLM trading agents, including FinMem's, come from one test window of a few months, a single back-test per stock and no transaction costs. A design that looks strong under such conditions may not survive a strict point-in-time, cost-aware, multi-regime test, so a new memory must be judged by a stricter standard.

The problem statement is therefore: \emph{Can a self-correcting agent memory that reasons about source credibility, redundancy, contradiction, related firms and its own past outcomes give an LLM trading agent better context and safer decisions than existing layered memory, when it is evaluated under a rigorous point-in-time protocol on recent data, and can it be deployed safely in an agentic trading application?}

\section{Motivation}
Retail and institutional users increasingly experiment with LLM agents for investing, yet published evidence is dominated by short back-tests whose results are hard to reproduce. Computationally, the memory is the component that decides \emph{which} evidence the model sees; poor memory produces redundant, stale or contradictory context, which wastes the context window and can bias the decision. Better memory is also useful outside trading: any agent that consumes a stream of documents of varying reliability (news monitoring, compliance, intelligence analysis) faces the same issues. Finally, an honest re-evaluation is valuable in itself: if a widely cited architecture does not transfer to new data, practitioners should know before deploying it.

\section{Objectives}
\begin{enumerate}
  \item Build a reproducible, point-in-time dataset of news, SEC filings and prices for FinMem's five stocks plus five additional large caps over about one year.
  \item Design a self-correcting agent memory with seven components (N1--N7) addressing credibility, redundancy, contradiction, query formation, importance, related firms and outcome feedback.
  \item Implement strong memory baselines on the same data: a faithful re-implementation of the layered FinMem memory (layered decay, importance, purge rules, access-counter promotion, top-$K$ per layer) and flat retrieval.
  \item Compare all systems against rule-based, machine-learning and LLM baselines under two evaluation protocols with transaction costs and appropriate statistical tests, and quantify each component through ablation.
  \item Test whether the conclusions depend on the model family by evaluating deep-learning (GRU, Transformer), reinforcement-learning (Q-learning) and pattern-matching models on exactly the same protocol, and analyse robustness (market exposure, stock set, market regime, statistical power).
  \item Deliver a working application in which the agent's retrieved evidence, recommendation and LLM decision can be inspected and executed on a paper-trading account, and in which the agent can trade autonomously only inside user-defined budget limits and permission levels.
  \item Make the application usable by non-specialists: plain-language explanations, budget-allocation advice, an algorithm board with honest track records, a Bangla interface and safe API-key management.
\end{enumerate}

\section{Methodology}
We downloaded 23{,}247 Benzinga articles through the Alpaca News API (the same source used by FinMem), 93 XBRL-based 10-Q/10-K summaries from SEC EDGAR and consolidated daily bars for ten stocks, SPY and QQQ. Each item becomes visible only after the first market close following its publication. All systems share the same FinBERT polarity scores and MiniLM sentence embeddings so that differences come only from the memory. Our memory and two baselines (flat retrieval and a faithful re-implementation of FinMem) are replayed day by day over the whole period; their retrieved evidence is turned into decisions by a common threshold rule, by gradient-boosting models and by three free LLM backbones. The primary evaluation is a rolling walk-forward over December~2025--October~2026 in which the decision threshold is re-selected every month from the previous 42 sessions; a fixed train/validation/test split is reported as a secondary protocol. Results are compared with Wilcoxon signed-rank tests (Holm-corrected), Diebold--Mariano tests, block-bootstrap confidence intervals and information coefficients. To check that the conclusions are not an artefact of the model family, a GRU, two Transformers, a Q-learning agent and a nearest-neighbour pattern model are added on the same protocol, and the results are decomposed into market exposure and alpha (CAPM), recomputed without Tesla, split by market regime and related to the statistical power of a one-year test. Finally, the research agent is packaged as an application with a safe-autonomy layer (budget guardrails, permission modes, approval queue, emergency stop) and analytics for non-specialist users.

\section{Project Outcome}
\begin{itemize}
  \item An open, point-in-time research dataset and a test-covered Python package (\testcount{} automated tests) that reproduces every number, table, figure and screenshot in this report.
  \item A self-correcting agent memory that is the best of the evaluated memory architectures on risk-adjusted return and drawdown, retrieves context about twice as fresh as the layered baseline's while merging re-reported stories and retiring contradicted facts, and improves the decisions of every LLM backbone tested, with an ablation showing which components matter.
  \item A rigorous evaluation framework for memory-based trading agents (point-in-time data, walk-forward protocol, costs, significance and power analysis), which shows honestly that the memory's advantage is not yet statistically significant, that the next-day predictive content of news memory is close to zero, and that the layered FinMem baseline does not transfer to 2025--26 data (Sharpe $-0.72$ on ten stocks, $-28.6\%$ on its own five stocks).
  \item A comparison with deep-learning, reinforcement-learning and pattern-matching models on the same protocol, together with CAPM, leave-one-out, regime and power analyses, which shows that the conclusions are not specific to one model family or one stock and quantifies how much data a significant result would need.
  \item A web application (``Agentic Trade'') that turns the research memory into an inspectable trading assistant connected to Alpaca paper trading, with a safe-autonomy layer (budget guardrails, four permission levels, approval queue, emergency stop), budget-allocation advice, analytics and an algorithm board for non-specialists, in English and Bangla.
\end{itemize}

\section{Organization of the Report}
Chapter~\ref{ch:back} reviews the background on financial text analysis, retrieval-augmented and memory-augmented LLM agents and evaluation methodology, and identifies the gap. Chapter~\ref{ch:design} presents the requirements, system design, the user interface, the detailed methodology (data, the baseline memories, our components N1--N7, decision policies, the safe-autonomy design and evaluation), the project plan and the task allocation. Chapter~\ref{ch:impl} describes the implementation, testing and the experimental results. Chapter~\ref{ch:sic} discusses the relevant standards, design constraints, cost and the complex-engineering-problem mapping. Chapter~\ref{ch:concl} concludes with limitations and future work. The appendices explain how to reproduce the results, list all hyperparameters and API endpoints, and give a short user guide.
